% !TEX program = lualatex
\documentclass[aspectratio=43,professionalfonts,handout]{beamer}
\usepackage{beamer_template}

\newcommand{\LectureNum}{14}
\newcommand{\LectureTitle}{スペクトル}
\newcommand{\TextPages}{p.\,100-103}
\newcommand{\CourseName}{応用数学}
\renewcommand{\TermName}{前期}

\begin{document}

\begin{frame}{\CourseName\quad 第\LectureNum 回「\LectureTitle」}
  {\small \TermName\quad \TeacherName\quad ／\quad 教科書 \TextPages}

  \textbf{目標}：スペクトルの概念を理解し、周期関数・非周期関数のスペクトルを求めることができる

\end{frame}

\begin{frame}{スペクトルとは}
  \begin{block}{}
  \begin{itemize}
    \item 関数を周波数成分に分解したもの
    \item 音・電気信号・画像処理への応用
    \item 周期関数 $ \to $ 線スペクトル
    \item 非周期関数 $ \to $ 連続スペクトル
  \end{itemize}
  \end{block}
\end{frame}

\begin{frame}{線スペクトル（周期関数）}
  \begin{block}{}
  \begin{itemize}
    \item フーリエ係数の絶対値をプロット
    \item 例2：矩形波のスペクトル
    \item $\omega = (2k-1)\pi/2$ にのみ成分が現れる
  \end{itemize}
  \end{block}
\end{frame}

\begin{frame}{連続スペクトル（非周期関数）}
  \begin{block}{}
  \begin{itemize}
    \item フーリエ変換 $F(u)$ がスペクトル密度
    \item 例3：$S_f(\omega) = 2sin \omega/(\pi\omega)$
    \item サンプリング定理との関係
  \end{itemize}
  \end{block}
\end{frame}

\begin{frame}{離散フーリエ変換（DFT）・FFT}
  \begin{block}{}
  \begin{itemize}
    \item デジタル信号処理の基礎
    \item FFT：高速アルゴリズム
    \item 画像処理・音声処理への応用
  \end{itemize}
  \end{block}
\end{frame}

\begin{frame}{演習（練習問題2）}
  \begin{exampleblock}{自分で解いてみよう}
  \begin{itemize}
    \item p.104 問1〜2より選題
  \end{itemize}
  \end{exampleblock}
\end{frame}

\begin{frame}{まとめ}
  \begin{itemize}
    \item 線スペクトルと連続スペクトルの違い
    \item 次回：前期総復習
  \end{itemize}
\end{frame}

\end{document}
